\section{Related Work}
\label{sec:related_work}

%\begin{table*}[!t]
%\caption{Comparison of existing checkpointing systems.
%}
%\label{tab:baseline_compare}
%\resizebox{\linewidth}{!}{
%\begin{tabular}{ccccc}
%    \toprule
%    \textbf{System} & \textbf{Multi-Parallel Strategy} & \textbf{Multi-Framework} & \textbf{Online Resharding} & \textbf{Scalabilty} \\
%    \midrule
%     CheckFreq~\cite{checkfreq} & \without & \without & \without & \without \\
%     Check-N-Run~\cite{check-n-run} & \without & \without & \without & \with \\
%     Gemini~\cite{gemini-sys} & \without & \without & \without & \without \\
%     JIT Checkpointing~\cite{jit-check} & \without & \without & \without & \without \\
%     DCP~\cite{torch-dcp} & \without & \without & \with & \without \\
%     MCP~\cite{megatron-dcp} & \without & \without & \with & \without \\
%     \sys{} & \with & \with & \with & \with \\
%    \bottomrule
%\end{tabular}}
%\end{table*}
%\vspace{-3mm}
%We summarize the features of different checkpointing systems in Table~\ref{tab:baseline_compare} and give a detailed discussion about existing works as follows: %\yhpeng{what does general purpose in the table mean? Consider using another term, e.g., large-scale?} 
% \yhpeng{The table can be deleted.}

% \yhpeng{we need more references here}

\noindent \textbf{Checkpointing frameworks.} %\yhpeng{discuss how PyTorch does checkpointing before DCP; also briefly discuss DeepSpeed UCP saying it is offline script}
Some industrial initiatives focus on developing checkpointing systems for deep learning. Prior to DCP~\cite{torch-dcp}, PyTorch provided \texttt{torch.save} and \texttt{torch.load} APIs for local checkpoint management but lacked native support for resharding. DCP~\cite{torch-dcp} introduces resharding capabilities for FSDP but does not support %hybrid
parallelism strategies like TP and PP. DeepSpeed-UCP~\cite{ucp} offers a unified format for DeepSpeed checkpoints and provides resharding capabilities through offline scripts. Megatron MCP~\cite{megatron-dcp} builds upon the workflow of DCP~\cite{torch-dcp} and extends available file storage formats to include options like Zarr~\cite{Zarr}. %However, 
All these frameworks %are constrained by their
provide limited support for various parallelisms and training frameworks, %(e.g., MCP only works for Megatron-LM~\cite{megatron}). 
% \zmf{the "limited support" and the reasons why "they cannot fulfill the demand" can be elaborated more in details.}
% Therefore, they cannot meet the multi-framework support requirement.
% Moreover, both DCP and MCP struggle with complex, irregular-sharded tensors in hybrid parallel training (see details in Sec.~\ref{sec:representation}), hindering their applicability in general scenarios.
%Additionally, when deployed for
and their I/O performance does not scale well in large-scale training.  %remains sub-optimal, and scalability becomes an issue. 
% Additionally, the centralized communication pattern used for process orchestration (See Sec.~\ref{sec:workflow}) in DCP~\cite{torch-dcp} and MCP~\cite{megatron-dcp} also presents potential stability risks. \sys{} address these limitations, enabling flexible multi-framework checkpointing support, as well as guaranteeing stability and efficiency at scale.

\noindent \textbf{Checkpointing optimizations.} Several %previous
works \cite{projectadam, deepfreeze} %\yhpeng{add citations of 2-3 papers here} 
 investigated reducing checkpointing costs from different perspectives.
Check-N-Run~\cite{check-n-run}, tailored for recommendation models, employs differential checkpointing to store only the modified portions of the model, alongside quantization %techniques
to reduce checkpoint size. % However, differential checkpointing is not directly applicable to the Pre-Training of LLMs, as LLMs typically exhibit dense updating patterns. The quantization technique is beyond the scope of this work, given that lossless checkpoints are imperative in our production settings.
CheckFreq~\cite{checkfreq} pipelines model states snapshot and saving operations with computation to reduce checkpoint stalls, and designs an online checkpoint frequency tuning algorithm to further reduce cost. 
Gemini~\cite{gemini-sys} advocates in-memory checkpointing with inter-machine backup for fast failure recovery and interleaves checkpoint communication traffic for state backup with training traffic, enabling frequent checkpointing in each iteration. JIT~\cite{jit-check} adopts just-in-time checkpointing for low-cost error recovery. 
%\mingji{Both CheckFreq~\cite{checkfreq} and Gemini~\cite{gemini-sys} focus on checkpoint efficiency during training.} 
%However, Gemini~\cite{gemini-sys}'s in-memory checkpointing solution demonstrates vulnerabilities in real-world large-scale GPU clusters, where multiple machines within the same backup group can experience simultaneous failures, leading to unacceptable catastrophic state loss. Other checkpointing techniques like Just-In-Time checkpointing  also suffer from the same reliability problem as Gemini~\cite{gemini-sys}. %\yhpeng{do not argue in this way regarding to Gemini and JIT. Maybe we can say these techniques are complementary and we can also use these techniques for better failure recovery} 
For industrial use cases, it is important to preserve checkpoints to persistent storage due to %the necessity of 
accessing checkpoints for various %steps for 
purposes like auto-evaluation~\cite{characterization}, hyper-parameter tuning, or even restarting training.
%Regarding the latter, after making configuration adjustments, the checkpoint with the highest evaluation scores is likely to be selected as the new starting point. 
Different from existing works, \sys{} champions both optimized I/O performance and flexible automatic online checkpoint resharding, supporting checkpointing for general LFM development.

\noindent \textbf{Checkpoint representation.} PyTorch's \texttt{torch.save/load} functionality relies on \texttt{pickle} for serialization and deserialization. %However, 
This format lacks crucial tensor shard metadata, such as global shape information, precluding automatic resharding. DCP \cite{torch-dcp} introduces a disaggregated format that separates metadata from tensor data. The metadata, which includes global shape and offset information, enables DCP's auto-resharding capabilities.
Built upon DCP's design, BCP incorporates necessary adaptations to support irregular tensor sharding. Array-based storage systems like Zarr \cite{zarrZarr} and TensorStore \cite{googleTensorStore} allow for saving tensors as individual arrays and support concurrent array reading and writing. MCP \cite{megatron-dcp} supports distributed checkpointing using the Zarr format.
For efficient and secure tensor storage, Safetensors \cite{huggingfaceSafetensors} is a file foramt for saving tensors safely and loading efficiently. For better compatibility with the Hugging Face open-source ecosystem, \sys{} incorporates functionality to export checkpoints in the Safetensors format.

\noindent \textbf{Storage systems for LFM.} Large-scale checkpointing demands robust storage systems. The LLaMA 3.1 report \cite{llama3.1} highlights Tectonic \cite{tectonic}, Meta's general-purpose file system, as the backbone for storing checkpoints during pre-training. One of the primary challenges faced by this system is managing high-frequency checkpoint writes. Similarly, FireFlyer \cite{fireflyer}, the AI-HPC system developed by DeepSeek \cite{deepseekai} for LFM training, incorporates the custom-built 3FS Distributed File System, akin to many parallel file systems such as BeeGFS \cite{herold2014introduction}, to handle checkpoint storage. %\sys{} leverages HDFS, implementing a series of performance optimizations to ensure scalability for large-scale checkpointing operations.

% \yhpeng{TODO: add related work for the above paragraphs}


% \textcolor{blue}{Mingji: A better categories could be: 1. Persistent Checkpoint system (online and offline resharding) 2. In memory fault checkpoint system }